\documentclass[12pt,a4paper]{article}
\usepackage[T5]{fontenc}
\usepackage[utf8]{inputenc}
\usepackage[vietnamese]{babel}
\usepackage{amsmath,amssymb}
\usepackage{geometry}
\usepackage{booktabs}
\usepackage{array}
\usepackage{tikz}
\usetikzlibrary{arrows.meta,positioning}
\usepackage{hyperref}
\geometry{margin=2.3cm}
\setlength{\parskip}{4pt}

\title{Tóm tắt kiến thức dự án\\\large Face Liveness Detection: phân biệt khuôn mặt thật (REAL) và giả (FAKE)}
\author{}
\date{}

\begin{document}
\maketitle

\section{Bài toán}
\textbf{Face liveness detection} (chống giả mạo khuôn mặt, \emph{face anti-spoofing}) kiểm tra khuôn mặt trước camera là
\begin{itemize}
  \item \textbf{REAL}: người thật đang đứng trước camera;
  \item \textbf{FAKE}: khuôn mặt giả, gọi là \emph{tấn công trình diện} (presentation attack): ảnh in, ảnh/video phát lại trên màn hình, mặt nạ.
\end{itemize}
Đây là bài toán \textbf{phân loại nhị phân}. Quy ước: lớp dương (positive) là REAL.

\section{Luồng xử lý một khung hình}
\begin{center}
\begin{tikzpicture}[node distance=0.55cm and 0.5cm, box/.style={draw, rounded corners, align=center, minimum height=1cm, font=\small}, >=Stealth]
  \node[box] (cam) {Camera\\(khung hình)};
  \node[box, right=of cam] (mp) {MediaPipe\\Face Landmarker};
  \node[box, above right=0.2cm and 0.6cm of mp] (st) {Mắt, quay đầu,\\ngẩng/cúi (hiển thị)};
  \node[box, below right=0.2cm and 0.6cm of mp] (mv) {MiniVision\\$p_{\text{real}}$};
  \node[box, right=of mv] (med) {Median\\15 khung};
  \node[box, right=of med] (dec) {$\ge T$: REAL\\$<T$: FAKE};
  \draw[->] (cam) -- (mp);
  \draw[->] (mp) |- (st);
  \draw[->] (mp) |- node[pos=0.75, below, font=\scriptsize]{hộp mặt} (mv);
  \draw[->] (mv) -- (med);
  \draw[->] (med) -- (dec);
\end{tikzpicture}
\end{center}

\section{Phát hiện khuôn mặt và điểm mốc}
\textbf{Google ML Kit} chỉ có SDK cho Android/iOS, \textbf{không có SDK Python}. Vì vậy dự án dùng \textbf{MediaPipe Face Landmarker} (cũng của Google). Với mỗi khuôn mặt, MediaPipe trả về:
\begin{itemize}
  \item 478 điểm mốc (landmark) $(u_i, v_i) \in [0,1]^2$; đổi ra pixel: $x_i = u_i W$, $y_i = v_i H$;
  \item các hệ số biểu cảm (\emph{blendshape}, mỗi hệ số từ 0 đến 1), trong đó có \texttt{eyeBlinkLeft}, \texttt{eyeBlinkRight};
  \item ma trận biến đổi khuôn mặt $4\times4$ (gồm phép xoay và tịnh tiến của đầu).
\end{itemize}
\textbf{Hộp mặt} lấy từ các điểm mốc:
\[
x_0 = \min_i x_i,\quad y_0 = \min_i y_i,\quad w = \max_i x_i - x_0,\quad h = \max_i y_i - y_0 .
\]

\section{Trạng thái mắt}
Độ nhắm mắt là trung bình hai blendshape:
\[
b = \frac{\texttt{eyeBlinkLeft} + \texttt{eyeBlinkRight}}{2} \in [0,1]
\qquad (0 = \text{mở},\ 1 = \text{nhắm}).
\]
\[
\text{Mắt} =
\begin{cases}
\text{CLOSED (nhắm)} & b \ge 0.5\\
\text{OPEN (mở)} & b < 0.5
\end{cases}
\]
Một lần \textbf{chớp mắt} là chuyển từ mở sang nhắm rồi mở lại trong thời gian ngắn.

\section{Tư thế đầu: quay trái/phải, ngẩng/cúi}
Lấy ma trận xoay $R$ ($3\times3$, góc trên trái của ma trận biến đổi), chia mỗi cột cho độ dài của nó để bỏ hệ số tỉ lệ. Các \textbf{góc Euler} (đơn vị độ):
\[
\text{yaw} = \operatorname{atan2}(R_{02}, R_{22}),\qquad
\text{pitch} = \arcsin(-R_{12}).
\]
\begin{itemize}
  \item \textbf{yaw}: quay ngang; $\text{yaw} > 0$ là quay sang phải ảnh.
  \item \textbf{pitch}: cúi/ngẩng; $\text{pitch} > 0$ là cúi xuống.
\end{itemize}
Dấu của hai góc đã được kiểm tra trên 400 ảnh: yaw tương quan $0.82$ với độ lệch ngang của mũi, pitch tương quan $0.81$ với độ lệch dọc của mũi.
\[
\text{Quay} =
\begin{cases}
\text{RIGHT} & \text{yaw} > 12^\circ\\
\text{LEFT} & \text{yaw} < -12^\circ\\
\text{CENTER} & \text{còn lại}
\end{cases}
\qquad
\text{Ngẩng/cúi} =
\begin{cases}
\text{DOWN (cúi)} & \text{pitch} > 10^\circ\\
\text{UP (ngẩng)} & \text{pitch} < -10^\circ\\
\text{CENTER} & \text{còn lại}
\end{cases}
\]
Các ngưỡng $0.5$, $12^\circ$, $10^\circ$ là \textbf{lựa chọn thiết kế}: dataset không có nhãn ``nhắm mắt'' hay ``quay đầu'' để học ngưỡng.

\textbf{Lưu ý quan trọng:} chớp mắt và xoay đầu \textbf{không chứng minh} người thật, vì một video phát lại người đang chớp mắt hay quay đầu vẫn vượt qua. Các tín hiệu này chỉ để hiển thị.

\section{MiniVision: chấm điểm người thật}
\textbf{MiniVision Silent-Face-Anti-Spoofing} dùng hai mạng CNN nhẹ, \textbf{MiniFASNetV2} và \textbf{MiniFASNetV1SE} (có khối \emph{Squeeze-and-Excitation}), huấn luyện sẵn. Mạng chỉ nhìn \textbf{một khung hình} và học kết cấu ảnh (vân màn hình, phản chiếu, mép ảnh in), nên người dùng không cần làm động tác (\emph{silent}).

\textbf{Cắt vùng mặt} với tỉ lệ mở rộng $s \in \{2.7,\ 4.0\}$ (mỗi mạng một tỉ lệ):
\[
k = \min\!\Big(\tfrac{H-1}{h}, \tfrac{W-1}{w}, s\Big),\quad
c_x = x_0 + \tfrac{w}{2},\quad c_y = y_0 + \tfrac{h}{2},
\]
\[
\text{vùng cắt} = \big[c_x - \tfrac{kw}{2},\ c_x + \tfrac{kw}{2}\big] \times \big[c_y - \tfrac{kh}{2},\ c_y + \tfrac{kh}{2}\big]\ \text{(cắt trong ảnh)},
\]
rồi resize về $80\times80$. Mỗi mạng cho 3 giá trị đầu ra $z_0, z_1, z_2$. Hàm \textbf{softmax} đổi chúng thành xác suất:
\[
p_j = \frac{e^{z_j}}{\sum_{m=0}^{2} e^{z_m}} .
\]
Lớp số 1 là ``người thật''. Lấy trung bình hai mạng:
\[
p_{\text{real}} = \frac{p_1^{(\text{V2})} + p_1^{(\text{V1SE})}}{2} \in [0,1].
\]

\section{Quyết định REAL/FAKE}
\[
\text{Kết luận} =
\begin{cases}
\text{REAL} & p_{\text{real}} \ge T\\
\text{FAKE} & p_{\text{real}} < T
\end{cases}
\]
Khi chạy realtime, để chống nhiễu từ một khung lẻ, hệ thống dùng \textbf{trung vị (median)} của 15 khung gần nhất,
$\tilde p = \operatorname{median}\big(p_{\text{real}}^{(t-14)}, \dots, p_{\text{real}}^{(t)}\big)$,
và chỉ kết luận khi đã có ít nhất 8 khung.

\section{Đánh giá và chọn ngưỡng}
\subsection{Ma trận nhầm lẫn (Confusion Matrix)}
\begin{center}
\begin{tabular}{lcc}
\toprule
 & Dự đoán FAKE & Dự đoán REAL \\
\midrule
Thực tế FAKE & TN & FP \\
Thực tế REAL & FN & TP \\
\bottomrule
\end{tabular}
\end{center}

\subsection{Các chỉ số}
\begin{align*}
\text{TPR} &= \frac{TP}{TP+FN} && \text{(tỉ lệ REAL được chấp nhận)}\\
\text{FPR} = \text{FAR} &= \frac{FP}{FP+TN} && \text{(FAKE bị nhận nhầm là REAL, nguy hiểm nhất)}\\
\text{FRR} &= \frac{FN}{FN+TP} = 1-\text{TPR} && \text{(REAL bị từ chối nhầm)}\\
\text{Accuracy} &= \frac{TP+TN}{TP+TN+FP+FN}
\end{align*}

\subsection{ROC và AUC}
\begin{itemize}
  \item \textbf{Đường ROC} là tập các điểm $\big(\text{FPR}(T), \text{TPR}(T)\big)$ khi cho ngưỡng $T$ chạy từ 1 về 0. Đường càng gần góc trên trái càng tốt.
  \item \textbf{AUC} là diện tích dưới đường ROC. AUC không phụ thuộc ngưỡng, và bằng xác suất một ảnh REAL ngẫu nhiên có điểm cao hơn một ảnh FAKE ngẫu nhiên:
  \[
  \text{AUC} = P\big(p_{\text{real}}^{\text{REAL}} > p_{\text{real}}^{\text{FAKE}}\big).
  \]
  $\text{AUC}=0.5$: đoán ngẫu nhiên; $\text{AUC}=1$: tách hoàn hảo.
\end{itemize}

\subsection{Chọn ngưỡng tốt nhất: chỉ số Youden}
\[
J(T) = \text{TPR}(T) - \text{FPR}(T) = 1 - \text{FRR}(T) - \text{FAR}(T),
\qquad
T^{*} = \arg\max_{T} J(T).
\]
Chọn $T^{*}$ tức là làm tổng hai loại lỗi FAR + FRR nhỏ nhất, coi hai loại lỗi quan trọng như nhau. Các giá trị $T$ được thử là chính điểm $p_{\text{real}}$ của từng ảnh.

\textbf{Quy trình đúng:} tìm $T^{*}$ trên \textbf{tập train}, rồi giữ nguyên để đo trên \textbf{tập test}. Nếu chọn ngưỡng trên chính tập test, kết quả sẽ lạc quan giả tạo.

\section{Dataset và kết quả}
\textbf{NUAA Photograph Imposter Database} (Tan và cộng sự, ECCV 2010): 15 người, ảnh $640\times480$ từ webcam. REAL là người thật, FAKE là ảnh in của người đó được cầm trước camera.

\begin{center}
\begin{tabular}{lrr}
\toprule
Tập & REAL & FAKE \\
\midrule
Train & 1\,743 & 1\,748 \\
Test & 3\,362 & 5\,761 \\
\bottomrule
\end{tabular}
\end{center}

\begin{center}
\begin{tabular}{lr}
\toprule
Chỉ số (trên tập test) & Giá trị \\
\midrule
Ngưỡng $T^{*}$ (Youden trên train) & 0.108 \\
Accuracy & 0.9993 \\
AUC & 1.000 \\
FAR & 0.10\% (6/5\,760) \\
FRR & 0\% (0/3\,362) \\
\bottomrule
\end{tabular}
\end{center}

Trên train, điểm REAL thấp nhất là $0.108$ và điểm FAKE cao nhất là $0.050$, nên Youden chọn $T^{*}=0.108$. Trên test, mọi ngưỡng trong khoảng $0.28$ đến $0.65$ đều không có lỗi nào.

\section{Hạn chế}
\begin{itemize}
  \item NUAA chỉ có tấn công bằng \textbf{ảnh in}, không có video phát lại hay màn hình. MiniVision tách hai lớp gần như hoàn hảo, nên ngưỡng $0.108$ \textbf{khá lỏng} khi dùng với camera thật; nên thử ngưỡng cao hơn, ví dụ $0.9$.
  \item Ngưỡng mắt và góc đầu chưa được kiểm chứng bằng dữ liệu có nhãn.
  \item Không rõ MiniVision có từng được huấn luyện trên NUAA hay không.
  \item Chưa thử realtime với người thật trước webcam trong lần chạy tự động.
\end{itemize}

\section*{Tài liệu tham khảo}
\begin{enumerate}
  \item MiniVision, \emph{Silent-Face-Anti-Spoofing}, \url{https://github.com/minivision-ai/Silent-Face-Anti-Spoofing} (Apache-2.0).
  \item Google, \emph{MediaPipe Face Landmarker}, \url{https://ai.google.dev/edge/mediapipe/solutions/vision/face_landmarker}.
  \item X. Tan, Y. Li, J. Liu, L. Jiang, \emph{Face Liveness Detection from a Single Image with Sparse Low Rank Bilinear Discriminative Model}, ECCV 2010.
  \item W. J. Youden, \emph{Index for rating diagnostic tests}, Cancer, 1950.
\end{enumerate}

\end{document}
